\documentclass[11pt]{article}

\usepackage[margin=1in]{geometry}
\usepackage[utf8]{inputenc}
\usepackage[T1]{fontenc}
\usepackage{amsmath,amssymb,amsthm}
\usepackage{graphicx}
\usepackage{xcolor}
\usepackage[colorlinks=true, citecolor=blue, linkcolor=black, urlcolor=blue]{hyperref}

\title{Rough notes}
\author{}
\date{\today}

\begin{document}
\maketitle

\section{The Models}

\begin{enumerate}
  \item \textbf{NK (Rotemberg).} Basic New Keynesian model with capital and
  Rotemberg quadratic price-adjustment costs. Rotemberg, ``Sticky Prices in the
  United States,'' \emph{JPE} (1982); textbook form in Gal\'i, \emph{Monetary
  Policy, Inflation, and the Business Cycle} (2015).

  \item \textbf{Calvo.} New Keynesian model with capital and Calvo staggered
  pricing. Calvo, ``Staggered Prices in a Utility-Maximizing Framework,''
  \emph{JME} (1983). The specific version solved here is Section~5 of Kahou,
  Fern\'andez-Villaverde, Gomez-Cram, Perla and Rosa, ``Spooky Boundaries at a
  Distance.''

  \item \textbf{NKFV.} Medium-scale New Keynesian model with stochastic
  volatility (unit-root technology trend, internal habits, intertemporal
  preference shock). Fern\'andez-Villaverde, Guerr\'on-Quintana, Kuester and
  Rubio-Ram\'irez, ``Fiscal Volatility Shocks and Economic Activity,''
  \emph{AER} 105(11) (2015).

  \item \textbf{Calvo + regimes.} The Calvo model augmented with a cost-push
  AR(1) shock and a two-state Markov supply regime in the firm's cost wedge.
  Nu\~no, Renner and Scheidegger, ``Monetary Policy with Supply Regimes.''
\end{enumerate}

\newpage
\section{NK (Rotemberg)}

\subsection{The equations}

\paragraph{Household.} The representative household maximizes
\begin{equation*}
\mathbb{E}_0 \sum_{t=0}^{\infty} \beta^t \left\{ \log c_t - \phi \frac{l_t^{1+\eta}}{1+\eta} \right\}
\end{equation*}
subject to
\begin{equation*}
c_t + k_t + \frac{b_{t+1}}{p_t} = w_t l_t + (1 + r_t - \delta) k_{t-1} + R_{t-1} \frac{b_t}{p_t} + T_t + F_t .
\end{equation*}
With $\lambda_t$ the multiplier on the budget constraint and $\Pi_t = p_t/p_{t-1}$, the first-order conditions are
\begin{align}
\frac{1}{c_t} &= \lambda_t, \\
\phi\, l_t^{\eta} &= \lambda_t w_t, \\
\lambda_t &= \beta\, \mathbb{E}_t \left[ \lambda_{t+1} \frac{R_t}{\Pi_{t+1}} \right], \\
\lambda_t &= \beta\, \mathbb{E}_t \left[ \lambda_{t+1} (1 + r_{t+1} - \delta) \right].
\end{align}

\paragraph{Final good producer.} The final good is $y_t = \left( \int_0^1 y_{it}^{\frac{\epsilon-1}{\epsilon}} di \right)^{\frac{\epsilon}{\epsilon-1}}$. The demand for variety $i$ and the price index are
\begin{equation}
y_{it} = \left( \frac{p_{it}}{p_t} \right)^{-\epsilon} y_t, \qquad
p_t = \left( \int_0^1 p_{it}^{1-\epsilon} di \right)^{\frac{1}{1-\epsilon}} .
\end{equation}

\paragraph{Intermediate good producers.} Firm $i$ produces $y_{it} = e^{z_t} k_{it-1}^{\alpha} l_{it}^{1-\alpha}$. Cost minimization gives the factor ratio and the real marginal cost, common across firms,
\begin{align}
\frac{k_{t-1}}{l_t} &= \frac{\alpha}{1-\alpha} \frac{w_t}{r_t}, \\
mc_t &= \left( \frac{1}{1-\alpha} \right)^{1-\alpha} \left( \frac{1}{\alpha} \right)^{\alpha} \frac{w_t^{1-\alpha} r_t^{\alpha}}{e^{z_t}} .
\end{align}
The firm chooses its price subject to a quadratic adjustment cost,
\begin{equation*}
\max_{\{p_{it}\}} \ \mathbb{E}_0 \sum_{t=0}^{\infty} \beta^t \lambda_t \left\{ \left( \frac{p_{it}}{p_t} \right)^{1-\epsilon} y_t - mc_t \left( \frac{p_{it}}{p_t} \right)^{-\epsilon} y_t - \frac{\psi}{2} \left( \frac{p_{it}}{\Phi_t\, p_{it-1}} - 1 \right)^2 y_t \right\},
\end{equation*}
where the indexed price target is
\begin{equation}
\Phi_t = \Pi^{1-\chi}\, \Pi_{t-1}^{\chi} .
\end{equation}
In the symmetric equilibrium the first-order condition is the New Keynesian Phillips curve
\begin{equation}
\psi \left( \frac{\Pi_t}{\Phi_t} - 1 \right) \frac{\Pi_t}{\Phi_t}
= (1 - \epsilon) + \epsilon\, mc_t
+ \psi \beta\, \mathbb{E}_t \left[ \frac{\lambda_{t+1}}{\lambda_t} \left( \frac{\Pi_{t+1}}{\Phi_{t+1}} - 1 \right) \frac{\Pi_{t+1}}{\Phi_{t+1}} \frac{y_{t+1}}{y_t} \right] .
\end{equation}

\paragraph{Government.} The central bank follows the Taylor rule
\begin{equation}
R_t = R \left( \frac{\Pi_t}{\Pi} \right)^{\gamma_\Pi} \left( \frac{y_t}{y_{t-1}} \right)^{\gamma_{\Delta y}} e^{\xi_t},
\end{equation}
where $R$ and $\Pi$ are the steady-state values. This is the output-growth rule, and it is the one the code uses, with $\gamma_\Pi = 2$ and $\gamma_{\Delta y} = 0.5$. The note \texttt{NK\_rotemberg.tex} uses the output-level rule instead, which replaces $y_t/y_{t-1}$ by $y_t/y$, where $y$ is the steady-state value of output.

\paragraph{Market clearing.}
\begin{align}
y_t &= e^{z_t} k_{t-1}^{\alpha} l_t^{1-\alpha}, \\
c_t + k_t + \frac{\psi}{2} \left( \frac{\Pi_t}{\Phi_t} - 1 \right)^2 y_t &= y_t + (1-\delta) k_{t-1} .
\end{align}

\paragraph{Shocks.}
\begin{align}
z_t &= \rho z_{t-1} + \sigma_z \varepsilon_{z,t}, \\
\xi_t &= \rho_m \xi_{t-1} + \sigma_m \varepsilon_{m,t},
\end{align}
with $\varepsilon_{z,t}, \varepsilon_{m,t} \sim \mathcal{N}(0,1)$.

\paragraph{State and policy.} The state is $s_t = (z_t, k_{t-1}, \xi_t, \Pi_{t-1}, y_{t-1})$. Lagged inflation $\Pi_{t-1}$ is a state because it enters the reference rate $\Phi_t$ of the adjustment cost, and $y_{t-1}$ is a state because it enters the Taylor rule. The policy is $(k_t, l_t, \Pi_t)$.

\subsection{Derivations}

\subsubsection{The household problem}

\paragraph{The terms.} In the objective, $\beta \in (0,1)$ is the discount factor and $\mathbb{E}_0$ is the expectation over the future shocks given the information at date 0. The term $\log c_t$ is the utility from consumption $c_t$. The term $\phi\, l_t^{1+\eta}/(1+\eta)$ is the disutility from hours worked $l_t$, where $\phi$ scales the disutility and $\eta$ is the inverse of the Frisch elasticity of labor supply. The left-hand side of the budget constraint is spending, all in units of the final good: consumption $c_t$, the capital $k_t$ carried into the next period, and the purchase of nominal bonds $b_{t+1}$, divided by the price level $p_t$ to convert them into goods. The right-hand side is resources. The term $w_t l_t$ is labor income at the real wage $w_t$. The term $(1 + r_t - \delta) k_{t-1}$ is the capital chosen in the previous period, which earns the real rental rate $r_t$ and loses the fraction $\delta$ to depreciation. The term $R_{t-1} b_t/p_t$ is the real value of the bonds bought in the previous period, which pay the gross nominal rate $R_{t-1}$ that was set at date $t-1$. The last two terms are the lump-sum transfers from the government $T_t$ and the profits of the intermediate firms $F_t$, which the household owns.

\paragraph{States and choices.} The household enters date $t$ with two individual state variables, the capital stock $k_{t-1}$ and the bond holdings $b_t$, which carry the predetermined rate $R_{t-1}$. Both were chosen at date $t-1$ and cannot be changed at date $t$. The choice variables at date $t$ are consumption $c_t$, hours $l_t$, capital $k_t$ and bonds $b_{t+1}$. The household takes as given the prices $w_t$, $r_t$, $p_t$ and $R_t$, the transfers $T_t$ and the profits $F_t$. These are functions of the aggregate state of the economy, which the household cannot affect: the shocks $z_t$ and $\xi_t$, the aggregate capital stock, lagged inflation $\Pi_{t-1}$ and lagged output $y_{t-1}$. In equilibrium the individual capital stock equals the aggregate capital stock. Bonds drop out of the resource constraint because the government adjusts the lump-sum transfers to service its debt, so the level of debt does not affect the allocation. This is why the equilibrium state $s_t$ contains $k_{t-1}$ but not $b_t$.

The household chooses $\{c_t, l_t, k_t, b_{t+1}\}_{t \ge 0}$ and takes prices, transfers $T_t$ and profits $F_t$ as given. The Lagrangian is
\begin{multline*}
\mathcal{L} = \mathbb{E}_0 \sum_{t=0}^{\infty} \beta^t \bigg\{ \log c_t - \phi \frac{l_t^{1+\eta}}{1+\eta} \\
+ \lambda_t \left[ w_t l_t + (1 + r_t - \delta) k_{t-1} + R_{t-1} \frac{b_t}{p_t} + T_t + F_t - c_t - k_t - \frac{b_{t+1}}{p_t} \right] \bigg\} .
\end{multline*}

\paragraph{Consumption.} The derivative with respect to $c_t$ is
\begin{equation*}
\frac{1}{c_t} - \lambda_t = 0 .
\end{equation*}

\paragraph{Labor.} The derivative with respect to $l_t$ is
\begin{equation*}
-\phi\, l_t^{\eta} + \lambda_t w_t = 0 .
\end{equation*}

\paragraph{Bonds.} The bond $b_{t+1}$ appears in the date $t$ constraint as a purchase and in the date $t+1$ constraint as a payoff. The derivative is
\begin{equation*}
-\frac{\lambda_t}{p_t} + \beta\, \mathbb{E}_t \left[ \lambda_{t+1} \frac{R_t}{p_{t+1}} \right] = 0 .
\end{equation*}
Multiply by $p_t$ and use $\Pi_{t+1} = p_{t+1}/p_t$ to get
\begin{equation*}
\lambda_t = \beta\, \mathbb{E}_t \left[ \lambda_{t+1} \frac{R_t}{\Pi_{t+1}} \right] .
\end{equation*}

\paragraph{Capital.} Capital $k_t$ is bought at date $t$ and pays $1 + r_{t+1} - \delta$ at date $t+1$. The derivative is
\begin{equation*}
-\lambda_t + \beta\, \mathbb{E}_t \left[ \lambda_{t+1} (1 + r_{t+1} - \delta) \right] = 0 .
\end{equation*}

\paragraph{Transversality.} The two Euler equations are necessary but not sufficient. The household's optimum also requires
\begin{equation*}
\lim_{T \to \infty} \mathbb{E}_0 \left[ \beta^T \lambda_T k_T \right] = 0, \qquad
\lim_{T \to \infty} \mathbb{E}_0 \left[ \beta^T \lambda_T \frac{b_{T+1}}{p_T} \right] = 0 .
\end{equation*}

\subsubsection{The final good producer}

The final good producer takes $p_t$ and $\{p_{it}\}$ as given and solves
\begin{equation*}
\max_{\{y_{it}\}} \ p_t \left( \int_0^1 y_{it}^{\frac{\epsilon-1}{\epsilon}} di \right)^{\frac{\epsilon}{\epsilon-1}} - \int_0^1 p_{it}\, y_{it}\, di .
\end{equation*}
The derivative of the aggregator with respect to $y_{it}$ is
\begin{equation*}
\frac{\partial y_t}{\partial y_{it}} = \left( \int_0^1 y_{jt}^{\frac{\epsilon-1}{\epsilon}} dj \right)^{\frac{1}{\epsilon-1}} y_{it}^{-\frac{1}{\epsilon}} = y_t^{\frac{1}{\epsilon}}\, y_{it}^{-\frac{1}{\epsilon}},
\end{equation*}
so the first-order condition $p_t\, y_t^{1/\epsilon}\, y_{it}^{-1/\epsilon} = p_{it}$ gives the demand function
\begin{equation*}
y_{it} = \left( \frac{p_{it}}{p_t} \right)^{-\epsilon} y_t .
\end{equation*}
The technology has constant returns, so profits are zero, $p_t y_t = \int_0^1 p_{it} y_{it}\, di$. Substitute the demand function,
\begin{equation*}
p_t y_t = \int_0^1 p_{it}^{1-\epsilon}\, p_t^{\epsilon}\, y_t\, di
\quad \Longrightarrow \quad
p_t = \left( \int_0^1 p_{it}^{1-\epsilon} di \right)^{\frac{1}{1-\epsilon}} .
\end{equation*}

\subsubsection{The intermediate firm: cost minimization}

Firm $i$ rents capital and labor in competitive markets. For a given output level $\bar y_{it}$ it solves
\begin{equation*}
\min_{k_{it-1},\, l_{it}} \ w_t l_{it} + r_t k_{it-1}
\quad \text{subject to} \quad e^{z_t} k_{it-1}^{\alpha} l_{it}^{1-\alpha} = \bar y_{it} .
\end{equation*}
Let $\mu_t$ be the multiplier on the constraint. The first-order conditions are
\begin{align*}
w_t &= \mu_t (1-\alpha)\, e^{z_t} k_{it-1}^{\alpha} l_{it}^{-\alpha}, \\
r_t &= \mu_t\, \alpha\, e^{z_t} k_{it-1}^{\alpha-1} l_{it}^{1-\alpha} .
\end{align*}
Divide the first by the second,
\begin{equation*}
\frac{w_t}{r_t} = \frac{1-\alpha}{\alpha} \frac{k_{it-1}}{l_{it}}
\quad \Longrightarrow \quad
\frac{k_{it-1}}{l_{it}} = \frac{\alpha}{1-\alpha} \frac{w_t}{r_t} .
\end{equation*}
The capital-labor ratio does not depend on $i$, so it equals the aggregate ratio $k_{t-1}/l_t$.

The multiplier $\mu_t$ is the cost of one more unit of output, that is, the real marginal cost $mc_t$. Solve the labor condition for $\mu_t$ and substitute the capital-labor ratio,
\begin{equation*}
mc_t = \frac{w_t}{(1-\alpha)\, e^{z_t}} \left( \frac{k_{it-1}}{l_{it}} \right)^{-\alpha}
= \frac{w_t}{(1-\alpha)\, e^{z_t}} \left( \frac{1-\alpha}{\alpha} \frac{r_t}{w_t} \right)^{\alpha}
= \left( \frac{1}{1-\alpha} \right)^{1-\alpha} \left( \frac{1}{\alpha} \right)^{\alpha} \frac{w_t^{1-\alpha} r_t^{\alpha}}{e^{z_t}} .
\end{equation*}
Marginal cost is the same for all firms and does not depend on the output level. With constant returns the total cost is $w_t l_{it} + r_t k_{it-1} = mc_t\, y_{it}$.

\subsubsection{The intermediate firm: pricing}

\paragraph{What the firm chooses.} Firm $i$ takes $p_t$, $y_t$, $mc_t$ and the demand function $y_{it} = (p_{it}/p_t)^{-\epsilon} y_t$ as given, with $\epsilon > 1$. It chooses its price $p_{it}$ to maximize profits.

\paragraph{Profits.} Real profits of firm $i$ at date $t$ are
\begin{equation*}
\frac{p_{it}}{p_t}\, y_{it} - mc_t\, y_{it} - \frac{\psi}{2} \left( \frac{p_{it}}{\Phi_t\, p_{it-1}} - 1 \right)^2 y_t .
\end{equation*}
The first term is the revenue. In dollars it is $p_{it}\, y_{it}$, and dividing by $p_t$ expresses it in units of the final good. The second term is the production cost $w_t l_{it} + r_t k_{it-1} = mc_t\, y_{it}$ from the cost minimization. The third term is the cost of changing the price. Substitute the demand function into the first two terms,
\begin{equation*}
\left( \frac{p_{it}}{p_t} \right)^{1-\epsilon} y_t - mc_t \left( \frac{p_{it}}{p_t} \right)^{-\epsilon} y_t - \frac{\psi}{2} \left( \frac{p_{it}}{\Phi_t\, p_{it-1}} - 1 \right)^2 y_t .
\end{equation*}

\paragraph{Where $\Phi_t$ comes from.} The adjustment cost is an assumption of the model. The familiar quadratic cost of changing a nominal price is $\frac{\psi}{2} (p_{it} - p_{it-1})^2$. Three changes lead from it to the term above.

First, the cost is on the percentage change of the price and not on the dollar change,
\begin{equation*}
\left( \frac{p_{it} - p_{it-1}}{p_{it-1}} \right)^2 = \left( \frac{p_{it}}{p_{it-1}} - 1 \right)^2 .
\end{equation*}
The dollar version is in dollars squared and quadruples when all prices double. The percentage version is a pure number.

Second, the cost is multiplied by aggregate output $y_t$, so that it is in units of the final good and stays proportional to the size of the economy.

Third, the reference price changes. So far the cost is zero only when $p_{it} = p_{it-1}$. Let $\Pi_t = p_t/p_{t-1}$ be gross inflation, that is, the growth rate of the price level between $t-1$ and $t$, and let $\Pi$ be its value in the steady state. If $\Pi \ne 1$, every firm has to raise its price at that rate just to keep its relative price, and it would pay the cost forever. The model replaces the reference price $p_{it-1}$ by $\Phi_t\, p_{it-1}$,
\begin{equation*}
\left( \frac{p_{it} - \Phi_t\, p_{it-1}}{\Phi_t\, p_{it-1}} \right)^2 = \left( \frac{p_{it}}{\Phi_t\, p_{it-1}} - 1 \right)^2 ,
\qquad
\Phi_t = \Pi^{1-\chi}\, \Pi_{t-1}^{\chi} , \quad \chi \in [0,1] .
\end{equation*}
The firm can raise its price at the rate $\Phi_t$ for free and pays only for the deviation from it. The rate $\Phi_t$ is a weighted geometric average of steady-state inflation $\Pi$ and last period's inflation $\Pi_{t-1}$. With $\chi = 0$ the free rate is $\Pi$, and with $\chi = 1$ it is $\Pi_{t-1}$. It depends only on aggregate variables that are known at date $t$, so the firm takes it as given.

\paragraph{Why the problem is dynamic.} The cost at date $t$ depends on the previous price $p_{it-1}$, so the cost at date $t+1$ depends on $p_{it}$. The firm cannot maximize date by date. It maximizes the value of the whole stream of profits.

\paragraph{Discounting.} The firm pays its profits to the household. One unit of the final good at date $t$ is worth $\lambda_t$ units of utility to the household, and utility at date $t$ is discounted by $\beta^t$. The firm therefore values the profits of date $t$ with the weight $\beta^t \lambda_t$ and, given the initial price $p_{i,-1}$, solves
\begin{equation*}
\max_{\{p_{it}\}} \ \mathbb{E}_0 \sum_{t=0}^{\infty} \beta^t \lambda_t \left\{ \left( \frac{p_{it}}{p_t} \right)^{1-\epsilon} y_t - mc_t \left( \frac{p_{it}}{p_t} \right)^{-\epsilon} y_t - \frac{\psi}{2} \left( \frac{p_{it}}{\Phi_t\, p_{it-1}} - 1 \right)^2 y_t \right\} .
\end{equation*}

\paragraph{The first-order condition.} The price $p_{it}$ appears in three places. It appears in the revenue and the production cost of date $t$, in the adjustment cost of date $t$ as the new price, and in the adjustment cost of date $t+1$ as the previous price. The derivative with respect to $p_{it}$ is
\begin{multline*}
\lambda_t \left\{ (1-\epsilon) \left( \frac{p_{it}}{p_t} \right)^{-\epsilon} \frac{y_t}{p_t}
+ \epsilon\, mc_t \left( \frac{p_{it}}{p_t} \right)^{-\epsilon-1} \frac{y_t}{p_t}
- \psi \left( \frac{p_{it}}{\Phi_t\, p_{it-1}} - 1 \right) \frac{y_t}{\Phi_t\, p_{it-1}} \right\} \\
+ \beta\, \mathbb{E}_t \left[ \lambda_{t+1}\, \psi \left( \frac{p_{it+1}}{\Phi_{t+1}\, p_{it}} - 1 \right) \frac{p_{it+1}}{\Phi_{t+1}\, p_{it}^2}\, y_{t+1} \right] = 0 .
\end{multline*}
The first two terms are the derivatives of the revenue and of the production cost. The third term is the extra adjustment cost paid today for a higher price. The last term is the adjustment cost saved tomorrow, because a higher price today means a smaller increase tomorrow. The benchmark $\Phi_{t+1}$ contains $\Pi_t$, which is aggregate inflation and not the firm's own price growth, so it is not differentiated.

\paragraph{Symmetric equilibrium.} All firms face the same problem. If they start from the same initial price they choose the same price at every date, and the price index gives $p_{it} = p_t$ for all $i$. Then $p_{it}/p_t = 1$ and $p_{it}/(\Phi_t\, p_{it-1}) = \Pi_t/\Phi_t$. Multiply the first-order condition by $p_t/(\lambda_t y_t)$,
\begin{equation*}
(1-\epsilon) + \epsilon\, mc_t - \psi \left( \frac{\Pi_t}{\Phi_t} - 1 \right) \frac{\Pi_t}{\Phi_t}
+ \psi \beta\, \mathbb{E}_t \left[ \frac{\lambda_{t+1}}{\lambda_t} \left( \frac{\Pi_{t+1}}{\Phi_{t+1}} - 1 \right) \frac{\Pi_{t+1}}{\Phi_{t+1}} \frac{y_{t+1}}{y_t} \right] = 0,
\end{equation*}
which is the Phillips curve of the previous subsection. Without the adjustment cost, $\psi = 0$, it reduces to $mc_t = (\epsilon-1)/\epsilon$, which says that the price is the markup $\epsilon/(\epsilon-1)$ over marginal cost. With the cost, the marginal cost can differ from $(\epsilon-1)/\epsilon$. When it is higher, firms want higher prices, and inflation rises above the benchmark $\Phi_t$ only as far as the adjustment costs of today and tomorrow allow.

\subsubsection{Aggregation and market clearing}

In the symmetric equilibrium all firms set the same price, so $y_{it} = y_t$, $k_{it-1} = k_{t-1}$ and $l_{it} = l_t$. There is no price dispersion, and aggregate output is
\begin{equation*}
y_t = e^{z_t} k_{t-1}^{\alpha} l_t^{1-\alpha} .
\end{equation*}
The household owns two kinds of firms. The final good producer has zero profits, so the profits $F_t$ in the household budget constraint are the profits of the intermediate firms. Real profits of intermediate firm $i$ are revenue minus factor payments minus the adjustment cost,
\begin{equation*}
F_{it} = \frac{p_{it}}{p_t}\, y_{it} - w_t l_{it} - r_t k_{it-1} - \frac{\psi}{2} \left( \frac{p_{it}}{\Phi_t\, p_{it-1}} - 1 \right)^2 y_t .
\end{equation*}
Integrate over the firms,
\begin{equation*}
F_t = \int_0^1 F_{it}\, di
= \frac{1}{p_t} \int_0^1 p_{it}\, y_{it}\, di - w_t \int_0^1 l_{it}\, di - r_t \int_0^1 k_{it-1}\, di - \frac{\psi}{2}\, y_t \int_0^1 \left( \frac{p_{it}}{\Phi_t\, p_{it-1}} - 1 \right)^2 di .
\end{equation*}
The first integral is $p_t y_t$ by the zero profit condition of the final good producer. The second and the third are the labor and capital market clearing conditions, $\int_0^1 l_{it}\, di = l_t$ and $\int_0^1 k_{it-1}\, di = k_{t-1}$. In the symmetric equilibrium $p_{it}/(\Phi_t\, p_{it-1}) = \Pi_t/\Phi_t$ for all $i$, so the last integral is $(\Pi_t/\Phi_t - 1)^2$. Aggregate profits are therefore output minus factor payments minus the adjustment cost,
\begin{equation*}
F_t = y_t - w_t l_t - r_t k_{t-1} - \frac{\psi}{2} \left( \frac{\Pi_t}{\Phi_t} - 1 \right)^2 y_t .
\end{equation*}
The government budget constraint is $T_t + R_{t-1} b_t/p_t = b_{t+1}/p_t$. Substitute both into the household budget constraint. The bond terms, the wage bill and the rental payments cancel, and what remains is the resource constraint
\begin{equation*}
c_t + k_t + \frac{\psi}{2} \left( \frac{\Pi_t}{\Phi_t} - 1 \right)^2 y_t = y_t + (1-\delta) k_{t-1} .
\end{equation*}

\subsubsection{Monetary policy}

The Taylor rule is not derived from an optimization problem. It is a behavioral rule that closes the model.

\subsubsection{Steady state}

Set all shocks to zero and all variables constant. Inflation equals the target, so $\Phi = \Pi$ and the adjustment cost is zero. The two Euler equations give
\begin{equation*}
R = \frac{\Pi}{\beta}, \qquad r = \frac{1}{\beta} - 1 + \delta .
\end{equation*}
The Phillips curve gives $mc = (\epsilon-1)/\epsilon$. Invert the marginal cost formula for the wage,
\begin{equation*}
w = \left[ (1-\alpha)^{1-\alpha} \alpha^{\alpha}\, \frac{mc}{r^{\alpha}} \right]^{\frac{1}{1-\alpha}} .
\end{equation*}
Normalize $l = 1$ and let $\phi$ adjust. Then
\begin{equation*}
k = \frac{\alpha}{1-\alpha} \frac{w}{r}, \qquad
y = k^{\alpha}, \qquad
c = y - \delta k, \qquad
\lambda = \frac{1}{c}, \qquad
\phi = \frac{w}{c} .
\end{equation*}

\end{document}
